\section{Introduction}\label{sec:intro}

Zero-knowledge virtual machines (zkVMs) enable a prover to execute arbitrary RISC-V programs and produce succinct cryptographic proofs of correct execution~\cite{risc0, sp1, jolt}. Production deployments already underpin blockchain scaling solutions, privacy-preserving applications, and off-chain computation platforms securing billions of dollars in digital assets~\cite{openvm, nexus, pico}. The security of these deployments hinges on a critical invariant: the zkVM's arithmetization---the polynomial constraints encoding the execution trace---must faithfully capture the RISC-V instruction semantics. When this alignment fails, a malicious prover can forge proofs for incorrect executions, undermining the entire trust model.

Detecting such \emph{semantic gaps} is extremely challenging. The constraint systems are large (hundreds of thousands of constraints), component interactions are subtle (e.g., the ALU must agree with the memory subsystem on the value written to a register), and bugs manifest only under specific \emph{boundary conditions}: rare combinations of operand values, instruction orderings, and execution states that exercise edge cases in the arithmetization. These bugs rarely surface during normal testing or standard security audits.

Existing approaches fall short. Static analysis tools such as \zkap~\cite{zkap} detect generic patterns (e.g., unconstrained signals) but lack zkVM-specific semantic structure, producing high false positive rates while missing actual bugs. Dynamic approaches such as \arguzz~\cite{arguzz} combine metamorphic testing with fault injection, but rely on generic mutation strategies unlikely to reach the boundary conditions where constraint bugs hide; in our experiments, \arguzz times out on the majority of benchmarks. Both lack a \emph{semantic model} of what constitutes a boundary condition and how to explore these conditions across heterogeneous backends.

Our key insight is that soundness-critical gaps in zkVMs arise at \emph{semantic boundary conditions} that become observable when heterogeneous backend traces are \emph{lifted into a common semantic space}. Although each zkVM represents its execution trace differently, the constraint-correctness concerns they share---immediate value decomposition, memory address alignment, control flow target binding, arithmetic overflow handling---can be captured by a small, well-defined set of semantic categories. Moreover, the bugs that these categories identify fall into two fundamentally distinct classes: \emph{divergences} (the executor computes incorrectly) and \emph{underconstraints} (the executor is correct but the constraint system is too permissive). Detecting both classes requires complementary mechanisms, a property we formalize as \emph{two-phase coverage} (\autoref{sec:formal-properties}). By naming semantic categories and using them as a structured feedback signal, a fuzzer can systematically explore the semantic dimensions where bugs cluster, rather than blindly mutating inputs.

Based on this insight, we introduce \tool, a semantics-guided differential fuzzing framework for discovering constraint bugs in RISC-V zkVMs. \tool makes two core contributions. First, it defines a \emph{cross-zkVM semantic model} that unifies six heterogeneous zkVM backends (OpenVM, SP1, RISC Zero, Jolt, Nexus, and Pico) into a common space of 17 observation types and 32 semantic categories, enabling cross-backend constraint-correctness reasoning without per-backend analysis (\autoref{sec:trace-ir}). Second, it builds a closed-loop \emph{boundary-driven fuzzing framework} on this shared abstraction: category signature novelty replaces edge coverage as the feedback signal, a UCB1 multi-armed bandit governs eight ISA-aware mutation strategies, and semantic witness injection probes underconstrained circuits at identified boundary points (\autoref{sec:fuzzing-framework}).

We evaluate \tool on all six backends against a ground-truth benchmark of 44 known vulnerabilities sourced from zero-day discovery, prior dynamic testing, security audits, and public issue trackers. \tool detects 36 of 44 bugs (81.8\% recall). On the four backends where Arguzz~\cite{arguzz} has full support, \tool detects 18 bugs to Arguzz's 11 (1.6$\times$); across all six backends, the improvement is 3.3$\times$ (36 vs.\ 11, 25.0\%). \tool additionally discovers 6 previously unknown zero-day vulnerabilities, including 3 critical soundness bugs in RISC Zero, OpenVM, and Pico that allow forging proofs for incorrect RISC-V execution. All zero-day vulnerabilities have been responsibly disclosed to the respective development teams.

In summary, we contribute:
\begin{itemize}
    \item A \textbf{cross-zkVM semantic model} of 17 observation types and 32 semantic categories that unifies six backends, with a formal two-phase coverage characterization (\autoref{sec:trace-ir}, \autoref{sec:formal-properties}).
    \item A \textbf{boundary-driven fuzzing framework} with category-signature feedback, bandit-guided ISA-aware mutation, and semantic witness injection (\autoref{sec:fuzzing-framework}).
    \item A \textbf{comprehensive evaluation} achieving 81.8\% recall on 44 known bugs ($3.3\times$ over the state-of-the-art) and 6 zero-day discoveries (\autoref{sec:eval}). We will open-source \tool upon acceptance.
\end{itemize}
